% !PRINTED 669
\weberpage{669}
%<<CONT
und durch Multiplication von (6) und (7) mit Rücksicht auf (4)

\weq{8}{e(r^n)\, e(r)^{-n} = \mathfrak{E}(r)^2.}

In unseren Betrachtungen über die Classenzahl wird sich noch der Satz ergeben:

\gesperrt{Eine Einheit $e(r)$ des reellen Körpers $H_m$, die mit allen ihren Conjugirten positiv ist, ist das Quadrat einer Einheit im Körper $H_m$.}

Die Formel (8) zeigt aber, dass der Einheit $\pm e(r)$ diese Eigenschaft zukommt, und dass daher $e(r)$ (da auch $-1 = i^2$ das Quadrat einer Einheit ist) das Quadrat einer Einheit des Körpers $\Omega_m$ ist.

Damit ist die Frage auf den Beweis der beiden erwähnten Sätze zurückgeführt, der sich als Schlussstein einer langen, aber an interessanten Beziehungen äusserst reichen Kette von Betrachtungen ergiebt, denen die nächsten Abschnitte gewidmet sein sollen.

\begin{center}
\rule{6em}{0.4pt}
\end{center}
